
\subsection{$A_\infty$ 同态之间的同伦}
\label{subsec:4-2-2}

从现在起，假定所有滤过 $A_\infty$ 代数及滤过 $A_\infty$ 同态都有隙。

\begin{definition}{4.2.35}\label{def:4-2-35}
令 $C_1,C_2$ 为滤过 $A_\infty$ 代数，令
$\mathfrak f,\mathfrak g\colon C_1\to C_2$ 为它们之间的滤过
$A_\infty$ 同态，再令 $\mathfrak C_2$ 为 $[0,1]\times C_2$ 的一个模型。
若存在滤过 $A_\infty$ 同态
$\mathfrak F\colon C_1\to\mathfrak C_2$，使得
\[
 \mathrm{Eval}_{s=0}\circ\mathfrak F=\mathfrak f,
 \qquad
 \mathrm{Eval}_{s=1}\circ\mathfrak F=\mathfrak g,
\]
就称 $\mathfrak f$ 经由 $\mathfrak C_2$ \emph{同伦于}
$\mathfrak g$，并记作
$\mathfrak f\sim_{\mathfrak C_2}\mathfrak g$。
称 $\mathfrak F$ 为 $\mathfrak f$ 与 $\mathfrak g$ 之间的
\emph{同伦}\liftIndex{I145}\liftIndex{I146}。
\end{definition}

\begin{lemma}{4.2.36}\label{lem:4-2-36}
关系 $\sim_{\mathfrak C_2}$ 不依赖于 $\mathfrak C_2$ 的选择。
\end{lemma}

\begin{proof}
令 $\mathfrak C'_2$ 为 $[0,1]\times C_2$ 的另一个模型。假定
$\mathfrak f\sim_{\mathfrak C_2}\mathfrak g$，并令 $\mathfrak F$ 如定义
4.2.35。把定理 4.2.34 用于恒等同态 $C_2\to C_2$，可得与
$\mathrm{Eval}$ 交换的同态
$\mathfrak{ID}\colon\mathfrak C_2\to\mathfrak C'_2$。置
$\mathfrak F'=\mathfrak{ID}\circ\mathfrak F$，则容易看出
\[
 \mathrm{Eval}_{s=0}\circ\mathfrak F'=\mathfrak f,
 \qquad
 \mathrm{Eval}_{s=1}\circ\mathfrak F'=\mathfrak g.
\]
故 $\mathfrak f\sim_{\mathfrak C'_2}\mathfrak g$，如所需。
\end{proof}

从现在起，以 $\sim$ 代替 $\sim_{\mathfrak C_2}$。

\begin{proposition}{4.2.37}\label{prop:4-2-37}
关系 $\sim$ 是等价关系。
\end{proposition}

\begin{proof}
$\mathfrak f\sim\mathfrak f$ 是显然的。

令 $\mathfrak f^{(m)}\colon C_1\to C_2$ 为滤过 $A_\infty$ 同态，并假定
$\mathfrak f^{(1)}\sim\mathfrak f^{(2)}$。取 $[0,1]\times C_2$ 的模型
$\mathfrak C_2$，以及从 $\mathfrak f^{(1)}$ 到 $\mathfrak f^{(2)}$ 的同伦
$\mathfrak F\colon C_1\to\mathfrak C_2$。如下考虑
$[0,1]\times C_2$ 的另一个模型 $\mathfrak C'_2$：作为滤过
$A_\infty$ 代数，置 $\mathfrak C'_2=\mathfrak C_2$；保持
$\mathrm{Incl}$ 不变，但交换 $\mathrm{Eval}_{s=0}$ 与
$\mathrm{Eval}_{s=1}$。也就是说，$\mathfrak C'_2$ 的
$\mathrm{Eval}_{s=0}$ 是 $\mathfrak C_2$ 的 $\mathrm{Eval}_{s=1}$，反之亦然。
于是，把 $\mathfrak F$ 看作从 $C_1$ 到 $\mathfrak C'_2$ 的滤过
$A_\infty$ 同态，它就是从 $\mathfrak f^{(2)}$ 到
$\mathfrak f^{(1)}$ 的同伦。因此
$\mathfrak f^{(2)}\sim\mathfrak f^{(1)}$。

再假定 $\mathfrak f^{(1)}\sim\mathfrak f^{(2)}$ 且
$\mathfrak f^{(2)}\sim\mathfrak f^{(3)}$。令
$\mathfrak F^{(m)}\colon C_1\to\mathfrak C_2$ 为从
$\mathfrak f^{(m)}$ 到 $\mathfrak f^{(m+1)}$ 的同伦。为了证明
$\mathfrak f^{(1)}\sim\mathfrak f^{(3)}$，定义滤过 $A_\infty$ 代数
$\mathfrak C'_2$ 如下：
{\small
\[
\begin{aligned}
 \mathfrak C'_2
  &={\bigl\{}(x,y)\in\mathfrak C_2\oplus\mathfrak C_2
    \mid \deg x=\deg y,\\
  &\hspace{7.7em}
    \mathrm{Eval}_{s=1}(x)=\mathrm{Eval}_{s=0}(y){\bigr\}},\\
 \mathfrak m'_k((x_1,y_1),\ldots,(x_k,y_k))
  &=\bigl(\mathfrak m_k(x_1,\ldots,x_k),
           \mathfrak m_k(y_1,\ldots,y_k)\bigr),\\
 \mathrm{Eval}'_{s=0}(x,y)&=\mathrm{Eval}_{s=0}(x),\\
 \mathrm{Eval}'_{s=1}(x,y)&=\mathrm{Eval}_{s=1}(y),\\
 \mathrm{Incl}'(x)&=(\mathrm{Incl}(x),\mathrm{Incl}(x)).
\end{aligned}
\]
}
容易看出，$\mathfrak C'_2$ 是 $[0,1]\times C_2$ 的一个模型。定义
$\mathfrak F^{(1)}\#\mathfrak F^{(2)}\colon C_1\to\mathfrak C'_2$ 为
\[
 (\mathfrak F^{(1)}\#\mathfrak F^{(2)})(z)
 =\bigl(\mathfrak F^{(1)}(z),\mathfrak F^{(2)}(z)\bigr).
\]
容易看出，$\mathfrak F^{(1)}\#\mathfrak F^{(2)}$ 是滤过
$A_\infty$ 同态，并且是从 $\mathfrak f^{(1)}$ 到
$\mathfrak f^{(3)}$ 的同伦。命题 4.2.37 得证。
\end{proof}

\begin{remark}{4.2.38}\label{rem:4-2-38}
在微分分次代数的情形，\textbf{[GrMo81]} 推论 10.7 证明了与命题
4.2.37 类似的结果。不过，在微分分次代数及微分分次同态的范畴中，只有再加
$C_1$ 为极小代数这一假设，该结果才成立。在（滤过）$A_\infty$ 代数的
范畴中工作，反而使问题简单了一些。
\end{remark}

现在简要比较我们的定义与文献中的其他定义。取
$\overline C_2^{[0,1]}$ 作为 $[0,1]\times\overline C_2$ 的模型，令
\[
 \overline{\mathfrak F}\colon\overline C_1
 \longrightarrow\overline C_2^{[0,1]}
\]
为 $A_\infty$ 同态，并满足
\[
 \overline{\mathfrak f}^{(m)}
 =\overline{\mathrm{Eval}}_{s=m}\circ\overline{\mathfrak F},
 \qquad m=0,1.
\]
回忆
\[
 \overline C_2^{[0,1]}
 =\overline{\mathfrak I}_0(\overline C_2)
  \oplus\overline{\mathfrak J}(\overline C_2)
  \oplus\overline{\mathfrak I}_1(\overline C_2).
\]
由端点评价条件，存在次数为 $-1$ 的同态\translatornote{原书把
$\overline{\mathfrak h}'_k$ 的次数写成 $+1$。但
$\overline{\mathfrak J}$ 的次数为 $+1$，而
$\overline{\mathfrak F}_k$ 保持次数，故
$\overline{\mathfrak h}'_k$ 必须具有次数 $-1$；这也与紧接的命题
4.2.40 一致，译文据此改正。}
\[
 \overline{\mathfrak h}'_k\colon
 B_k\overline C_1[1]\longrightarrow\overline C_2[1],
\]
使得
\[
 \overline{\mathfrak F}_k
 =\bigl(\overline{\mathfrak I}_0\circ
        \overline{\mathfrak f}^{(0)}_k\bigr)
  \oplus\bigl(\overline{\mathfrak J}\circ
        \overline{\mathfrak h}'_k\bigr)
  \oplus\bigl(\overline{\mathfrak I}_1\circ
        \overline{\mathfrak f}^{(1)}_k\bigr).
\]
置
\[
 \overline{\mathfrak h}_k(\mathbf x)
 =(-1)^{\deg'\mathbf x}\overline{\mathfrak h}'_k(\mathbf x).
\]

由 $\overline C_2^{[0,1]}$ 上 $A_\infty$ 结构的定义可得，
$\overline{\mathfrak F}$ 是 $A_\infty$ 同态，当且仅当下式成立：
{\small
\begin{equation}
\begin{aligned}
&\sum_c(-1)^{\deg'\mathbf x_c^{(3;1)}}
 \overline{\mathfrak m}\Bigl(
  \widehat{\overline{\mathfrak f}}{}^{(0)}
    (\mathbf x_c^{(3;1)}),
  \overline{\mathfrak h}(\mathbf x_c^{(3;2)}),
  \widehat{\overline{\mathfrak f}}{}^{(1)}
    (\mathbf x_c^{(3;3)})\Bigr)\\
&\qquad=\overline{\mathfrak f}^{(0)}(\mathbf x)
 -\overline{\mathfrak f}^{(1)}(\mathbf x)
 -\sum_c(-1)^{\deg'\mathbf x_c^{(3;1)}}
  \overline{\mathfrak h}\bigl(
   \mathbf x_c^{(3;1)},
   \overline{\mathfrak m}(\mathbf x_c^{(3;2)}),
   \mathbf x_c^{(3;3)}\bigr).
\end{aligned}
\tag{4.2.39}\label{eq:4-2-39}
\end{equation}
}
（式 \eqref{eq:4-2-39} 是从
$(\widehat{\overline{\mathfrak F}}\circ\widehat d
-\widehat d\circ\widehat{\overline{\mathfrak F}})(\mathbf x)=0$
中取 $\overline{\mathfrak J}(\overline C_2)$ 分量所得。）于是已经证明：

\begin{proposition}{4.2.40}\label{prop:4-2-40}
$A_\infty$ 同态
$\overline{\mathfrak f}^{(0)}\colon\overline C_1\to\overline C_2$
同伦于
$\overline{\mathfrak f}^{(1)}\colon\overline C_1\to\overline C_2$，当且仅当
存在次数为 $-1$ 的同态族
\[
 \overline{\mathfrak h}_k\colon
 B_k\overline C_1[1]\longrightarrow\overline C_2[1],
\]
满足 \eqref{eq:4-2-39}。滤过版本也成立。
\end{proposition}

式 \eqref{eq:4-2-39} 正是多篇文献采用的 $A_\infty$ 同态之间同伦的
定义，例如 \textbf{[Smi00]} 第 4.1 节；本书预印本
\textbf{[FOOO00]} 的第 A5 节也采用了它。

第一署名作者在 \textbf{[Fuk02II]} 中研究了 $A_\infty$ 范畴的同伦等价。
只有一个对象的 $A_\infty$ 范畴就是 $A_\infty$ 代数，而两个单对象
$A_\infty$ 范畴之间的 $A_\infty$ 函子就是 $A_\infty$ 同态。应用
\textbf{[Fuk02II]} 定义 8.5 中两个 $A_\infty$ 函子同伦的定义，容易看出
它也等价于存在满足 \eqref{eq:4-2-39} 的 $\overline{\mathfrak h}_k$。
事实上，两个 $A_\infty$ 函子同伦的定义要求存在满足适当性质的自然变换；
上面的同态 $\overline{\mathfrak h}_k$ 正是这样的自然变换。

下面写出使用模型
$\overline{\operatorname{Poly}}([0,1],\overline C_2)$ 时，两个滤过
$A_\infty$ 同态同伦的条件。$A_\infty$ 同态
\[
 \overline{\mathfrak h}\colon\overline C_1
 \longrightarrow\overline{\operatorname{Poly}}([0,1],\overline C_2)
\]
可写成
$\overline{\mathfrak h}=(\overline{\mathfrak f}(s),
\overline{\mathfrak g}(s))$，其中
\[
 \overline{\mathfrak f}(s),\overline{\mathfrak g}(s)
 \colon B\overline C_1[1]\longrightarrow\overline C_2[1][s].
\]
这里 $\overline{\mathfrak f}(s)$ 保持次数，而
$\overline{\mathfrak g}(s)$ 的次数为 $-1$。
\translatornote{原书声称 $\overline{\mathfrak f}(s)$ 与
$\overline{\mathfrak g}(s)$ 都保持次数。按定义 4.2.9，区间模型的第二分量
比第一分量低一次，故 $\overline{\mathfrak g}(s)$ 必须具有次数 $-1$；
如此式 (4.2.41.2) 两边才次数齐次。}
$\overline{\mathfrak h}$ 为 $A_\infty$ 同态的条件可写成
\begin{equation}
 \overline{\mathfrak f}(s_0)\colon\overline C_1\longrightarrow
 \overline C_2\text{ 对每个 }s_0\in[0,1]
 \text{ 都是 }A_\infty\text{ 同态},
\tag{4.2.41.1}\label{eq:4-2-41-1}
\end{equation}
以及
{\small
\begin{equation}
 \frac{\partial\overline{\mathfrak f}(s)}{\partial s}
 =\overline{\mathfrak m}\circ
  \bigl(\widehat{\overline{\mathfrak f}}(s)
   \otimes\widehat{\overline{\mathfrak g}}(s)
   \otimes\widehat{\overline{\mathfrak f}}(s)\bigr)
  \circ\Delta^2
 -\overline{\mathfrak g}(s)\circ\widehat d.
\tag{4.2.41.2}\label{eq:4-2-41-2}
\end{equation}
}
因为后文不用这些条件，这里不检验 (4.2.41)。

现在转而研究同伦的基本性质。由引理 4.2.36，同伦关系不依赖于
$[0,1]\times C$ 模型的选择。因此，从现在起，我们不指明具体选择，而统一用
花体字母 $\mathfrak C$ 表示 $[0,1]\times C$ 的模型；只有当这种独立性并非
显然时，才另行检验。

\begin{definition}{4.2.42}\label{def:4-2-42}
若存在 $A_\infty$ 同态
$\overline{\mathfrak g}\colon\overline C_2\to\overline C_1$，使
$\overline{\mathfrak f}\circ\overline{\mathfrak g}$ 与
$\overline{\mathfrak g}\circ\overline{\mathfrak f}$ 分别同伦于恒等同态，
就称 $A_\infty$ 同态
$\overline{\mathfrak f}\colon\overline C_1\to\overline C_2$ 为
\emph{同伦等价}\liftIndex{I147}。

若两个 $A_\infty$ 代数之间存在同伦等价，就称它们彼此同伦等价。
\end{definition}

有隙滤过 $A_\infty$ 代数之间的同伦等价也以同样方式定义。

为了证明 $A_\infty$ 代数的同伦等价是等价关系，需要下面的引理。

\begin{lemma}{4.2.43}\label{lem:4-2-43}
令
$\mathfrak f^{(m)},
 \mathfrak f^{(m)\prime}\colon
 C_m\to C_{m+1}$ 为 $A_\infty$ 同态。若
$\mathfrak f^{(m)}\sim
 \mathfrak f^{(m)\prime}$，则
\[
 \mathfrak f^{(2)}\circ\mathfrak f^{(1)}
 \sim
 \mathfrak f^{(2)\prime}\circ
 \mathfrak f^{(1)\prime}.
\]
滤过 $A_\infty$ 同态也有同样的结论。
\translatornote{原书本引理首句和末句都写“滤过 $A_\infty$ 同态”，字面近乎
重复；经 400 dpi 复核（2026-09-23），源书通篇并无横线记号，译文照录源书
原样，见疑误 A-0197（2026-09-23 更正）。}
\end{lemma}

\begin{proof}
容易看出
$\mathfrak f^{(2)}\circ\mathfrak f^{(1)}
 \sim\mathfrak f^{(2)\prime}\circ\mathfrak f^{(1)}$，故可以假定
$\mathfrak f^{(2)}=\mathfrak f^{(2)\prime}$。把定理 4.2.34 用于
$\mathfrak f^{(2)}$，得到 $A_\infty$ 同态
$\mathfrak F^{(2)}\colon\mathfrak C_2\to\mathfrak C_3$，满足
\[
 \mathrm{Eval}_{s=s_0}\circ\mathfrak F^{(2)}
 =\mathfrak f^{(2)}\circ\mathrm{Eval}_{s=s_0}.
\]
令 $\mathfrak F\colon C_1\to\mathfrak C_2$ 为
$\mathfrak f^{(1)}$ 与 $\mathfrak f^{(1)\prime}$ 之间的同伦，则
$\mathfrak F^{(2)}\circ\mathfrak F$ 是从
$\mathfrak f^{(2)}\circ\mathfrak f^{(1)}$ 到
$\mathfrak f^{(2)}\circ\mathfrak f^{(1)\prime}$ 的同伦。
\end{proof}

下面的推论是直接的。

\begin{corollary}{4.2.44}\label{cor:4-2-44}
\begin{enumerate}[label=\textup{(\arabic*)},leftmargin=2.8em]
 \item 同伦等价的复合仍是同伦等价。
 \item 同伦等价是等价关系。
\end{enumerate}
\end{corollary}

下面的定理将在第 4.5 节证明。回忆：无滤过 $A_\infty$ 代数之间的
$A_\infty$ 同态，若诱导上链同伦等价，就称为弱同伦等价；满足
$\overline{\mathfrak m}_0=0$ 的滤过 $A_\infty$ 代数之间的滤过
$A_\infty$ 同态 $\mathfrak f\colon C_1\to C_2$，若其经典约化
$\overline{\mathfrak f}\colon\overline C_1\to\overline C_2$ 是弱同伦等价，
就称 $\mathfrak f$ 为弱同伦等价（定义 3.2.10、3.2.33）。若
$\mathfrak f_0=0$，则称滤过 $A_\infty$ 同态 $\mathfrak f$ 为严格同态
（定义 3.2.29）。

\begin{theorem}{4.2.45}\label{thm:4-2-45}
（$A_\infty$ 代数的 Whitehead 定理）\liftIndex{I151}
\begin{enumerate}[label=\textup{(\arabic*)},leftmargin=2.8em]
 \item $A_\infty$ 代数的弱同伦等价是同伦等价。
 \item 有隙滤过 $A_\infty$ 代数之间的有隙弱同伦等价是同伦等价。严格弱
 同伦等价的同伦逆可以取为严格同态。
\end{enumerate}
\end{theorem}

文献中有一个类似结果：对任意弱同伦等价
$\overline{\mathfrak f}\colon\overline C_1\to\overline C_2$，存在
$A_\infty$ 同态
$\overline{\mathfrak g}\colon\overline C_2\to\overline C_1$，它在
$\overline{\mathfrak m}_1$-上同调上诱导逆映射。定理 4.2.45 (1) 比这一
陈述稍强，因为一个在 $\overline{\mathfrak m}_1$-上同调上诱导恒等映射的
$A_\infty$ 同态，未必同伦于恒等同态。

类似地，也可以定义（滤过）$A_\infty$ 双模的（滤过）$A_\infty$ 双模同态
之间的同伦；第 5 章第 5.2 节将讨论它。

\subsection{含幺或同伦含幺情形}
\label{subsec:4-2-3}

本小节专门研究含幺或同伦含幺（滤过）$A_\infty$ 同态的同伦；相关定义见
第 3.3.2 小节。本小节主要固定若干记号并导出后文要用的基本事实。这里虽只
明确讨论滤过版本，但所有相应讨论都不作改变地适用于无滤过情形。我们仍假定
所有滤过 $A_\infty$ 代数与滤过 $A_\infty$ 同态都有隙。

令 $C_1,C_2$ 为两个含幺滤过 $A_\infty$ 代数。取
$[0,1]\times C_i$ 的模型 $\mathfrak C_i$；可假定它们也含幺。

\begin{definition}{4.2.46}\label{def:4-2-46}
设 $\mathfrak f,\mathfrak g\colon C_1\to C_2$ 是含幺滤过
$A_\infty$ 同态。若存在含幺滤过 $A_\infty$ 同态
$\mathfrak h\colon C_1\to\mathfrak C_2$，使
\[
 \mathrm{Eval}_{s=0}\circ\mathfrak h=\mathfrak f,
 \qquad
 \mathrm{Eval}_{s=1}\circ\mathfrak h=\mathfrak g,
\]
就称 $\mathfrak f$ 与 $\mathfrak g$ \emph{含幺同伦}，并记作
$\mathfrak f\sim\mathfrak g$（或
$\mathfrak f\sim_u\mathfrak g$）\liftIndex{I150}。称 $\mathfrak h$ 为
$\mathfrak f$ 与 $\mathfrak g$ 之间的含幺同伦。

对严格滤过 $A_\infty$ 代数之间的同态，以同样方式定义严格含幺同伦，并以
$\sim_{su}$ 表示。

有时简称为 u-同伦、su-同伦，或称两个同态 u-同伦、su-同伦，以分别代替
含幺同伦、严格含幺同伦及其相应关系。
\end{definition}

\begin{lemma}{4.2.47}\label{lem:4-2-47}
关系 $\sim_u$ 不依赖于模型 $\mathfrak C_2$ 的选择，并且是等价关系。若
$\mathfrak f\sim_u\mathfrak f'$，则
\[
 \mathfrak f\circ\mathfrak g\sim_u\mathfrak f'\circ\mathfrak g,
 \qquad
 \mathfrak g'\circ\mathfrak f\sim_u\mathfrak g'\circ\mathfrak f'.
\]
严格版本也一样成立。
\end{lemma}

证明与引理 4.2.36、命题 4.2.37 及引理 4.2.43 的证明相同。现在可以按显然的
方式定义\emph{含幺同伦等价}（或 u-同伦等价），也可类似定义严格同伦等价
\liftIndex{I149}。

\begin{theorem}{4.2.45a}\label{thm:4-2-45a}
定理 4.2.45 的含幺版本成立。此外，若
$\mathfrak f\colon C_1\to C_2$ 是含幺滤过 $A_\infty$ 同态且为同伦等价，
则它是 u-同伦等价。严格含幺版本也成立。
\end{theorem}

注意，后半部分由前半部分推出。前半部分将在第 4.5 节证明。

令 $C_1,C_2$ 为两个同伦含幺滤过 $A_\infty$ 代数。按 (3.3.1) 取
$C_i^+$，并取 $[0,1]\times C_2^+$ 的含幺模型
$\mathfrak C_2^+$。
\translatornote{原书把 $C_i^+$ 的定义误指向第 8.1 节；第 8.1 节讨论无标记
圆盘模空间的定向，而这里使用的正是 (3.3.1) 所定义的含幺扩张，译文据此
改正交叉引用。}

\begin{remark}{4.2.48}\label{rem:4-2-48}
可以证明，存在 $[0,1]\times C_2$ 的同伦含幺模型，并且它在类似于定理
4.2.34 的意义下唯一。于是也可以用这个模型定义同伦含幺同伦。不过，采用
这里的方法就无需证明这一点。
\end{remark}

\begin{definition}{4.2.49}\label{def:4-2-49}
设 $\mathfrak f,\mathfrak g\colon C_1\to C_2$ 是同伦含幺滤过
$A_\infty$ 同态。若存在含幺滤过 $A_\infty$ 同态
\[
 \mathfrak h\colon C_1^+\longrightarrow\mathfrak C_2^+,
\]
使
\[
 \mathrm{Eval}_{s=0}\circ\mathfrak h=\mathfrak f^+,
 \qquad
 \mathrm{Eval}_{s=1}\circ\mathfrak h=\mathfrak g^+,
\]
就称 $\mathfrak f$ 与 $\mathfrak g$ \emph{同伦含幺地同伦}，并记作
$\mathfrak f\sim\mathfrak g$（或
$\mathfrak f\sim_{hu}\mathfrak g$）\liftIndex{I148}。这里
$\mathfrak f^+,\mathfrak g^+\colon C_1^+\to C_2^+$ 如定义 3.3.13。
称 $\mathfrak h$ 为 $\mathfrak f$ 与 $\mathfrak g$ 之间的同伦含幺同伦。

有时简称说两个同态“hu-同伦”，并把连接它们的同伦简称为“hu-同伦”；
严格版本以同样方式定义，相应简称为“shu-同伦”。
\end{definition}

\begin{lemma}{4.2.50}\label{lem:4-2-50}
关系 $\sim_{hu}$ 不依赖于模型 $\mathfrak C_2$ 的选择，并且是等价关系。若
$\mathfrak f\sim_{hu}\mathfrak f'$，则
\[
 \mathfrak f\circ\mathfrak g\sim_{hu}\mathfrak f'\circ\mathfrak g,
 \qquad
 \mathfrak g'\circ\mathfrak f\sim_{hu}\mathfrak g'\circ\mathfrak f'.
\]
对 $\sim_{shu}$ 也有同样的结论。
\end{lemma}

证明与引理 4.2.47 相同。现在可以按显然的方式定义同伦含幺同伦等价
（hu-同伦等价）。

\begin{theorem}{4.2.45b}\label{thm:4-2-45b}
定理 4.2.45 的同伦含幺版本成立。此外，若
$\mathfrak f\colon C_1\to C_2$ 是同伦含幺滤过 $A_\infty$ 同态且为同伦
等价，则它是 hu-同伦等价。严格版本及严格同伦含幺版本也成立。
\end{theorem}

证明将在第 4.5 节给出。继续之前，先提醒一个容易混淆、虽属技术性却很微妙的
问题。

\begin{remark}{4.2.51}\label{rem:4-2-51}
我们把同伦含幺滤过 $A_\infty$ 同态
$\mathfrak f\colon C_1\to C_2$ 定义为含幺滤过 $A_\infty$ 同态
$\mathfrak f^+\colon C_1^+\to C_2^+$。这意味着，对给定的
$\mathfrak f\colon C_1\to C_2$，$\mathfrak f^+$ 可以有多种选择；而且，
具有同一 $\mathfrak f$ 的不同 $\mathfrak f^+$ 被视为不同的同伦含幺滤过
$A_\infty$ 同态。作者不知道，对 $\mathfrak f^+$ 的另一个选择
$\mathfrak f^{+\prime}\colon C_1^+\to C_2^+$ 是否总与
$\mathfrak f^+$ 含幺同伦。不过，容易证明
$\mathfrak f^{+\prime}$ 与 $\mathfrak f^+$ 普通同伦。

把含幺滤过 $A_\infty$ 同态 $\mathfrak f$ 看作同伦含幺滤过
$A_\infty$ 同态时，我们取 $\mathfrak f^+$ 的典范选择，使
\[
 \mathfrak f^+(\cdots\otimes\mathbf f\otimes\cdots)=0.
\]
\end{remark}

依注 4.2.51 指定的选择，可以定义函子
\[
 \mathcal I_{\mathrm{unit}}\colon
 [\text{含幺滤过 }A_\infty\text{ 代数}]
 \longrightarrow
 [\text{同伦含幺滤过 }A_\infty\text{ 代数}].
\]
另一方面，$C\mapsto C^+$ 定义函子
\[
 \mathcal I_+\colon
 [\text{同伦含幺滤过 }A_\infty\text{ 代数}]
 \longrightarrow
 [\text{含幺滤过 }A_\infty\text{ 代数}].
\]

\begin{proposition}{4.2.52}\label{prop:4-2-52}
函子 $\mathcal I_{\mathrm{unit}}$ 与 $\mathcal I_+$ 诱导如下范畴等价：
{\small
\[
 \begin{gathered}
 [\text{含幺滤过 }A_\infty\text{ 代数；}\
 \text{含幺滤过 }A_\infty\text{ 同态的 u-同伦类}]\\[-0.2em]
 \cong\\[-0.2em]
 [\text{同伦含幺滤过 }A_\infty\text{ 代数；}\
 \text{同伦含幺滤过 }A_\infty\text{ 同态的 hu-同伦类}].
 \end{gathered}
\]
}
严格版本也成立。
\end{proposition}

\begin{proof}
先证明下面的引理。

\begin{lemma}{4.2.53}\label{lem:4-2-53}
设 $\mathfrak f,\mathfrak g\colon C\to C'$ 是含幺滤过
$A_\infty$ 代数之间的两个含幺滤过 $A_\infty$ 同态。它们同伦含幺地同伦，
当且仅当它们含幺同伦。
\end{lemma}

证明要用下面的子引理。

\begin{statement}{子引理}{4.2.54}\label{sublem:4-2-54}
若 $C$ 是含幺滤过 $A_\infty$ 代数，则存在严格含幺滤过
$A_\infty$ 同态
\[
 \mathfrak{ret}\colon C^+\longrightarrow C,
\]
使复合 $\mathfrak{ret}\circ i\colon C\to C$ 为恒等同态；这里
$i\colon C\to C^+$ 是典范嵌入。
\end{statement}

\begin{proof}[子引理 4.2.54 的证明]
置
\[
 \mathfrak{ret}_1(\mathbf e^+)=\mathbf e,
 \qquad
 \mathfrak{ret}_1(\mathbf f)=0,
 \qquad
 \mathfrak{ret}_k=0\quad(k\ne1),
\]
并在 $C\subset C^+$ 上置 $\mathfrak{ret}_1=\mathrm{id}$。容易看出，
$\mathfrak{ret}$ 具有所需性质。
\end{proof}

注意，$i$ 不含幺。

\begin{proof}[引理 4.2.53 的证明]
应用定理 4.2.45a，分别为 $\mathfrak{ret}$ 与 $\mathfrak{ret}'$ 选择含幺
同伦逆
\[
 \mathfrak{ret}^{-1}\colon C\longrightarrow C^+,
 \qquad
 \mathfrak{ret}'^{-1}\colon C'\longrightarrow C'{}^+.
\]
先假定 $\mathfrak f\sim_{hu}\mathfrak g$，则
$\mathfrak f^+\sim_u\mathfrak g^+$，因而
{\small
\begin{align*}
 \mathfrak f
 &\sim_u\mathfrak f\circ\mathfrak{ret}\circ\mathfrak{ret}^{-1}
  =\mathfrak{ret}'\circ\mathfrak f^+\circ\mathfrak{ret}^{-1}\\
 &\sim_u\mathfrak{ret}'\circ\mathfrak g^+\circ\mathfrak{ret}^{-1}
  =\mathfrak g\circ\mathfrak{ret}\circ\mathfrak{ret}^{-1}
  \sim_u\mathfrak g.
\end{align*}
}
这里，等式
$\mathfrak f\circ\mathfrak{ret}=\mathfrak{ret}'\circ\mathfrak f^+$
来自注 4.2.51 中对 $\mathfrak f^+$ 的选择以及上面对
$\mathfrak{ret}$ 的定义。

再假定 $\mathfrak f\sim_u\mathfrak g$，则
{\small
\begin{align*}
 \mathfrak f^+
 &\sim_u\mathfrak{ret}'^{-1}\circ\mathfrak{ret}'\circ\mathfrak f^+
  =\mathfrak{ret}'^{-1}\circ\mathfrak f\circ\mathfrak{ret}\\
 &\sim_u\mathfrak{ret}'^{-1}\circ\mathfrak g\circ\mathfrak{ret}
  =\mathfrak{ret}'^{-1}\circ\mathfrak{ret}'\circ\mathfrak g^+
  \sim_u\mathfrak g^+.
\end{align*}
}
因此 $\mathfrak f\sim_{hu}\mathfrak g$。
\end{proof}

引理 4.2.53 表明，$\mathcal I_{\mathrm{unit}}$ 在同伦范畴上诱导单射函子。
另一方面，容易看出 $\mathcal I_+$ 在同伦范畴之间诱导函子。
所以，为完成命题 4.2.52 的证明，只需证明下面的引理。事实上，引理
4.2.55 表明 $\mathcal I_{\mathrm{unit}}\circ\mathcal I_+$ 等价于恒等函子。

\begin{lemma}{4.2.55}\label{lem:4-2-55}
若 $C$ 是同伦含幺滤过 $A_\infty$ 代数，则 $C$ 与 $C^+$ hu-同伦等价。
\end{lemma}

\begin{proof}
如下把 $i\colon C\to C^+$ 延拓成含幺滤过 $A_\infty$ 同态
$i^+\colon C^+\to C^{++}$。写成
\[
 \begin{aligned}
 C^+&=C\oplus\Lambda_{0,\mathrm{nov}}\mathbf e^+
          \oplus\Lambda_{0,\mathrm{nov}}\mathbf f,\\
 C^{++}&=C^+\oplus\Lambda_{0,\mathrm{nov}}\mathbf e^{++}
          \oplus\Lambda_{0,\mathrm{nov}}\mathbf f^+,
 \end{aligned}
\]
并置
\[
 i_1^+(x)=x\quad(x\in C),
 \qquad
 i_1^+(\mathbf e^+)=\mathbf e^{++},
 \qquad
 i_1^+(\mathbf f)=\mathbf f^++\mathbf f,
\]
以及 $i_k^+=0$（$k\ne1$）。容易检验 $i^+$ 是含幺滤过
$A_\infty$ 同态；注意 $C^+$ 的幺元为 $\mathbf e^+$，而 $C^{++}$ 的幺元为
$\mathbf e^{++}$。因此 $i$ 是同伦含幺同态。由定理 4.2.45b，$i$ 是
hu-同伦等价。
\end{proof}

命题 4.2.52 得证。
\end{proof}

利用命题 4.2.52，可以把同伦含幺滤过 $A_\infty$ 代数的大多数代数讨论
约化为含幺情形。因此，只有在绝对必要时，我们才专门讨论同伦含幺滤过
$A_\infty$ 代数。
